\section{Methods}
\label{sec:methods}

All thirteen pipelines share the same normalized text, the same Kiwi morpheme tokens
(\cref{sec:preprocessing}) and the same training subsets. Every data-dependent
transform---TF-IDF vocabulary and IDF weights, the neural vocabulary---is fitted on
the current training subset only. Because the two families use different
representations, the comparison is between \emph{complete pipelines}.

\paragraph{TF-IDF family.}
Morpheme unigrams and bigrams, at most 10{,}000 features, sublinear term frequency,
$\ell_2$ normalisation, sparse matrices; all four run on the CPU. Four classifiers:
logistic regression (\textbf{LR}) and a linear SVM (\textbf{SVM}) from scikit-learn~\cite{pedregosa2011scikit},
Random Forest~\cite{breiman2001random} (\textbf{RF}) and
XGBoost~\cite{chen2016xgboost} (\textbf{XGB}, histogram trees). LR uses the
\texttt{liblinear} solver with an $\ell_2$ penalty, SVM is a linear SVM with squared hinge loss
(\texttt{LinearSVC}), RF averages fully grown or depth-limited trees over random feature
subsets, and XGB grows boosted trees with the logistic loss, row and column subsampling, and a number
of rounds fixed by early stopping during tuning.

\subsection{Neural architectures}
\label{sec:arch}

\paragraph{Shared template.}
All nine neural classifiers follow the same template (\cref{fig:arch}) and differ only in
the backbone. A review is a sequence of Kiwi morpheme ids $x_1,\dots,x_\ell$, truncated
to $L=\maxSeqLenKiwi{}$ and right-padded with id~0. The vocabulary holds at most
20{,}000 morphemes of the current training subset with frequency $\ge2$, plus
\texttt{<pad>} and \texttt{<unk>}. Each id is mapped to a 128-dimensional embedding
$\mathbf{e}_t=\mathbf{E}x_t$ learned from scratch (no pre-trained vectors; the padding row is
fixed at zero). The backbone turns $\mathbf{e}_{1:\ell}$ into hidden states
$\mathbf{h}_{1:\ell}$, which are averaged over the \emph{valid} positions only,
\begin{equation}
  \bar{\mathbf{h}} = \frac{1}{\ell}\sum_{t=1}^{\ell}\mathbf{h}_t ,
  \qquad
  \hat{\mathbf{y}} = \operatorname{softmax}\!\bigl(\mathbf{W}_o\,\operatorname{Dropout}(\bar{\mathbf{h}}) + \mathbf{b}_o\bigr) \in \mathbb{R}^2 ,
  \label{eq:pool}
\end{equation}
and a linear layer produces the two class logits. Masked mean pooling is used by every
backbone, so differences between models come from how $\mathbf{h}_t$ is computed, not
from the readout. Dropout with the tuned rate $p$ is applied to the backbone input and
to $\bar{\mathbf{h}}$. With $|V|\approx20{,}000$ the embedding table alone has
$\approx2.56$M parameters for every model; we therefore match the \emph{non-embedding}
parameters across architectures (below).

\begin{figure}[t]
  \centering
  \begin{tikzpicture}[
      node distance=4mm and 5mm,
      box/.style={draw, rounded corners=2pt, minimum height=7mm, align=center, font=\small, inner sep=3pt},
      bb/.style={box, fill=black!6, text width=33mm, minimum height=6mm, font=\footnotesize},
      arr/.style={-{Stealth[length=2mm]}, thick}]
    \node[box] (in) {Kiwi morphemes\\$x_1,\dots,x_\ell$};
    \node[box, right=of in] (emb) {Embedding\\$128$-d, from scratch};
    \node[bb, right=8mm of emb, yshift=11mm] (mlp) {MLP (identity)};
    \node[bb, below=1.5mm of mlp] (rnn) {(Bi)RNN / (Bi)LSTM /\\ (Bi)GRU};
    \node[bb, below=1.5mm of rnn] (trf) {Transformer $\times2$};
    \node[bb, below=1.5mm of trf] (mam) {Mamba $\times2$};
    \node[draw, dashed, rounded corners=3pt, fit=(mlp)(rnn)(trf)(mam), inner sep=2.5pt,
          label={[font=\scriptsize]above:backbone $\mathbf{h}_{1:\ell}$}] (back) {};
    \node[box, right=8mm of back] (pool) {Masked\\mean pool};
    \node[box, right=of pool] (head) {Dropout +\\Linear $\to2$};
    \draw[arr] (in) -- (emb);
    \draw[arr] (emb.east) -- (back.west |- emb.east);
    \draw[arr] (back.east |- pool.west) -- (pool);
    \draw[arr] (pool) -- (head);
  \end{tikzpicture}
  \caption{Shared template of the nine neural classifiers. Only the backbone differs; the
  input, the 128-d embedding, masked mean pooling over valid tokens and the linear head are
  identical. For the MLP the backbone is the identity and the two-layer MLP is applied after
  pooling (deep averaging network).}
  \label{fig:arch}
\end{figure}

\paragraph{MLP (deep averaging network).}
Following DAN~\cite{iyyer2015deep}, the embeddings are pooled first and the pooled vector
passes through two hidden layers of width $w$:
$\mathbf{z}_0=\bar{\mathbf{e}}$, $\mathbf{z}_k=\operatorname{Dropout}(\operatorname{ReLU}(\mathbf{W}_k\mathbf{z}_{k-1}+\mathbf{b}_k))$
for $k=1,2$, followed by the linear head. It ignores word order completely and is the
neural counterpart of a bag-of-words model.

\paragraph{Recurrent networks (RNN, LSTM, GRU and their bidirectional versions).}
Each recurrent model has one layer. The \emph{unidirectional} RNN, LSTM and GRU read the
sequence left-to-right only, $\mathbf{h}_t=\overrightarrow{\mathbf{h}}_t\in\mathbb{R}^{w}$, so
the state at position $t$ has seen only the prefix $x_{1:t}$. The \emph{bidirectional}
BiRNN, BiLSTM and BiGRU add a second, independent layer that reads right-to-left~\cite{schuster1997bidirectional};
$w$ is then the hidden size per direction and
$\mathbf{h}_t=[\overrightarrow{\mathbf{h}}_t;\overleftarrow{\mathbf{h}}_t]\in\mathbb{R}^{2w}$, so
every position sees the whole review. Comparing each pair isolates the value of
right-to-left context at equal parameter count. Sequences are packed, so padding never
enters the recurrence. Within each direction setting the three variants differ only in
the cell (written below for the forward direction). The \textbf{RNN}/\textbf{BiRNN} use the plain Elman cell~\cite{elman1990finding},
$\overrightarrow{\mathbf{h}}_t=\tanh(\mathbf{W}\mathbf{e}_t+\mathbf{U}\overrightarrow{\mathbf{h}}_{t-1}+\mathbf{b})$,
which has no gates and is the baseline for what gating adds.
The \textbf{LSTM}/\textbf{BiLSTM}~\cite{hochreiter1997long} add input, forget and output gates and a
separate memory cell,
$\mathbf{c}_t=\mathbf{f}_t\odot\mathbf{c}_{t-1}+\mathbf{i}_t\odot\tanh(\mathbf{W}_c\mathbf{e}_t+\mathbf{U}_c\mathbf{h}_{t-1}+\mathbf{b}_c)$,
$\mathbf{h}_t=\mathbf{o}_t\odot\tanh(\mathbf{c}_t)$, with each gate
$\sigma(\mathbf{W}_g\mathbf{e}_t+\mathbf{U}_g\mathbf{h}_{t-1}+\mathbf{b}_g)$.
The \textbf{GRU}/\textbf{BiGRU}~\cite{cho2014learning} merge memory and state and use an update
gate $\mathbf{u}_t$ and a reset gate $\mathbf{r}_t$:
$\tilde{\mathbf{h}}_t=\tanh(\mathbf{W}\mathbf{e}_t+\mathbf{U}(\mathbf{r}_t\odot\mathbf{h}_{t-1})+\mathbf{b})$,
$\mathbf{h}_t=(1-\mathbf{u}_t)\odot\mathbf{h}_{t-1}+\mathbf{u}_t\odot\tilde{\mathbf{h}}_t$.
Per direction, an RNN, GRU and LSTM layer has one, three and four weight blocks of
size $w(128+w+2)$ respectively, so at equal budget the plain RNN is the widest, and a
unidirectional model is wider than its bidirectional twin, which splits the budget over
two directions (\cref{tab:arch}).

\paragraph{Transformer encoder.}
The embedding is projected to $d=w$ dimensions and a learned positional embedding
$\mathbf{p}_t$ ($L$ positions) is added: $\mathbf{u}_t=\mathbf{W}_p\mathbf{e}_t+\mathbf{p}_t$.
Two pre-norm encoder layers~\cite{vaswani2017attention} follow, each
$\mathbf{a}=\mathbf{u}+\operatorname{MHSA}(\operatorname{LN}(\mathbf{u}))$,
$\mathbf{u}'=\mathbf{a}+\mathbf{W}_2\operatorname{GELU}(\mathbf{W}_1\operatorname{LN}(\mathbf{a}))$,
with four attention heads, a feed-forward width of $2d$, dropout $p$ inside both
sub-layers, and a key-padding mask so that attention ignores padding. A final
LayerNorm~\cite{ba2016layer} precedes pooling. Positional embeddings are counted as
embedding parameters.

\paragraph{Mamba.}
The embedding is projected to $d=w$ and passed through two pre-norm residual blocks
$\mathbf{x}\leftarrow\mathbf{x}+\operatorname{Mixer}(\operatorname{RMSNorm}(\mathbf{x}))$~\cite{zhang2019root}.
The mixer is the selective state-space layer of \citet{gu2023mamba} with expansion 2
($d_\text{in}=2d$), a causal depthwise convolution of width 4 and state size $N=16$:
$[\tilde{\mathbf{x}},\mathbf{z}]=\mathbf{W}_\text{in}\mathbf{x}$,
$\tilde{\mathbf{x}}\leftarrow\operatorname{SiLU}(\operatorname{conv}(\tilde{\mathbf{x}}))$, and for every
channel
\begin{equation}
  \mathbf{s}_t=\exp(\Delta_t\mathbf{A})\odot\mathbf{s}_{t-1}+\Delta_t\mathbf{B}_t\tilde{x}_t,\qquad
  y_t=\mathbf{C}_t^{\top}\mathbf{s}_t+D\,\tilde{x}_t ,
\end{equation}
where the step size $\Delta_t=\operatorname{softplus}(\cdot)$ (rank $\lceil d/16\rceil$) and
$\mathbf{B}_t,\mathbf{C}_t\in\mathbb{R}^{N}$ are computed from $\tilde{\mathbf{x}}_t$: this
input dependence (``selection'') lets the model decide per token what to keep in its
state. The output is $\mathbf{W}_\text{out}(\mathbf{y}\odot\operatorname{SiLU}(\mathbf{z}))$.
The scan is causal, so right padding cannot affect valid positions; a final RMSNorm
precedes pooling. Mamba has two interchangeable implementations with identical
parameters: the fused CUDA kernels of \texttt{mamba-ssm} and a pure-PyTorch reference
selective scan (\texttt{torch\_ref}); all reported Mamba results use \mambaBackend{}, and CPU
latency always uses \texttt{torch\_ref}.

\paragraph{Parameter matching.}
Each architecture has three width levels S/M/L. For every architecture and level we
search the width $w$ whose non-embedding parameter count is closest to 50k, 200k and
500k respectively; all resulting models lie within $\pm3\%$ of the target
(\cref{tab:arch}). The width level is a tuned hyperparameter (\cref{sec:methods}), so
each dataset may select a different level.

\begin{table}[t]
  \centering
  \caption{Neural architectures. ``Width'' is the size that is varied to meet the
  parameter budget; S/M/L give the width and, in parentheses, the resulting number of
  non-embedding parameters (target in the header). All models additionally have a
  $|V|\times128$ embedding table ($\approx$2.56M parameters at $|V|\approx20$k) and a
  linear output layer.}
  \label{tab:arch}
  % AUTO-GENERATED by ksent.analysis.make_tables (do not edit by hand).
\begin{adjustbox}{max width=\linewidth}
\begin{tabular}{l>{\raggedright\arraybackslash}p{5.2cm}>{\raggedright\arraybackslash}p{2.8cm}lrrr}
\toprule
Method & Backbone & Fixed structure & Width = & S (50.0k) & M (200.0k) & L (500.0k) \\
\midrule
Emb + MLP & masked mean $\to$ 2$\times$[Linear, ReLU, dropout] & 2 hidden layers & hidden units per layer & 167 (49.9k) & 386 (199.9k) & 644 (499.7k) \\
Emb + RNN & 1-layer Elman RNN ($\tanh$), masked mean & 1 layer, 1 dir. & hidden units & 167 (49.9k) & 386 (199.9k) & 644 (499.7k) \\
Emb + LSTM & 1-layer LSTM, masked mean & 1 layer, 1 dir. & hidden units & 64 (49.8k) & 168 (200.6k) & 294 (499.2k) \\
Emb + GRU & 1-layer GRU, masked mean & 1 layer, 1 dir. & hidden units & 79 (49.7k) & 201 (200.0k) & 348 (499.7k) \\
Emb + BiRNN & 1-layer bidirectional Elman RNN ($\tanh$), masked mean & 1 layer, 2 dir. & hidden units per direction & 105 (49.8k) & 257 (199.9k) & 438 (499.3k) \\
Emb + BiLSTM & 1-layer bidirectional LSTM, masked mean & 1 layer, 2 dir. & hidden units per direction & 37 (49.6k) & 106 (200.6k) & 193 (499.5k) \\
Emb + BiGRU & 1-layer bidirectional GRU, masked mean & 1 layer, 2 dir. & hidden units per direction & 47 (50.1k) & 129 (201.0k) & 231 (501.3k) \\
Transformer & Linear 128$\to d$ + learned positions, pre-norm encoder, LayerNorm, masked mean & 2 layers, 4 heads, FFN $2d$, GELU & $d_\text{model}$ & 52 (51.3k) & 108 (203.4k) & 172 (500.0k) \\
Mamba & Linear 128$\to d$, pre-norm residual Mamba blocks, RMSNorm, masked mean & 2 blocks, $d_\text{state}$=16, conv 4, expand 2 & $d_\text{model}$ & 50 (49.3k) & 112 (196.4k) & 186 (498.9k) \\
\bottomrule
\end{tabular}
\end{adjustbox}

\end{table}

\paragraph{Fixed neural training rules.}
AdamW~\cite{loshchilov2019decoupled} (weight decay 0.01), batch size 64, at most 30
epochs, early stopping with patience 5 on validation macro-F1 with the best
checkpoint restored, gradient clipping at 1.0, cross-entropy loss and fp32
precision---identical for every size.

\paragraph{Tuning.}
For each dataset and method we evaluate exactly ten predefined configurations
(\cref{tab:a2defs}) at $n=2{,}000$ with tuning seed 0 and select the one with the
highest validation macro-F1 (ties: lower fit time). The search covers
regularisation for linear models, tree structure and sampling for RF, depth,
learning rate and sampling for XGBoost (whose early-stopped number of rounds is
frozen), and learning rate, dropout and width for neural models. The selected
configuration is frozen for all sizes and seeds; the validation macro-F1 of every tuning
configuration is listed in \cref{tab:a2}.

\begin{sidewaystable}
  \centering
  \caption{Method configurations. ``Selected'' is the frozen configuration per
  dataset (RF: trees, depth, features, minimum leaf size, class weight; XGB: depth,
  learning rate, subsampling, rounds; neural: width level, learning rate, dropout).
  ``Size'' is the non-embedding parameter count (neural), the number of non-zero
  coefficients (LR/SVM) or total tree nodes (RF/XGB) at $n=8{,}000$.
  ``Tune'' is the summed fit time of the ten tuning runs.}
  \label{tab:methods}
  % AUTO-GENERATED by ksent.analysis.make_tables (do not edit by hand).
\begin{adjustbox}{max width=\linewidth}
\begin{tabular}{lllllrrlrrlrr}
\toprule
Method & Family & Representation & Search space & \multicolumn{3}{c}{NSMC} & \multicolumn{3}{c}{Shopping} & \multicolumn{3}{c}{Steam} \\
\cmidrule(lr){5-7}\cmidrule(lr){8-10}\cmidrule(lr){11-13}
 & & & & Selected & Size & Tune (s) & Selected & Size & Tune (s) & Selected & Size & Tune (s) \\
\midrule
TF-IDF + LR & TF-IDF & TF-IDF 1--2-gram & $C\in\{10^{-2},\dots,300\}$ & -- & -- & -- & -- & -- & -- & -- & -- & -- \\
TF-IDF + SVM & TF-IDF & TF-IDF 1--2-gram & $C\in\{10^{-2},\dots,300\}$ & -- & -- & -- & -- & -- & -- & -- & -- & -- \\
TF-IDF + RF & TF-IDF & TF-IDF 1--2-gram & trees, depth, feat., leaf, weight (10/32) & -- & -- & -- & -- & -- & -- & -- & -- & -- \\
TF-IDF + XGBoost & TF-IDF & TF-IDF 1--2-gram & depth, $\eta$, sampling (10/18); rounds by ES & -- & -- & -- & -- & -- & -- & -- & -- & -- \\
\midrule
Emb + MLP & Neural & Emb.\ mean pool & lr, dropout, width S/M/L (10/27) & -- & -- & -- & -- & -- & -- & -- & -- & -- \\
Emb + LSTM & Neural & Emb.\ + BiLSTM & lr, dropout, width S/M/L (10/27) & -- & -- & -- & -- & -- & -- & -- & -- & -- \\
Emb + GRU & Neural & Emb.\ + BiGRU & lr, dropout, width S/M/L (10/27) & -- & -- & -- & -- & -- & -- & -- & -- & -- \\
Transformer & Neural & Emb.\ + 2-layer encoder & lr, dropout, width S/M/L (10/27) & -- & -- & -- & -- & -- & -- & -- & -- & -- \\
Mamba & Neural & Emb.\ + 2 Mamba blocks & lr, dropout, width S/M/L (10/27) & -- & -- & -- & -- & -- & -- & -- & -- & -- \\
\bottomrule
\end{tabular}
\end{adjustbox}

\end{sidewaystable}
